\providecommand{\itescastudentname}{Martin Jonathan de la Cruz}
\providecommand{\itescastudentshort}{M. J. de la Cruz}
\providecommand{\itescastudentid}{26130503}
\providecommand{\itescafaculty}{Instituto Tecnologico Superior de Cajeme}
\providecommand{\itescafacultyname}{Maestria en Gestion Administrativa}
\providecommand{\itescacoursename}{Estadistica Aplicada para la Administracion}
\providecommand{\itescacoursecode}{MGA-ESTADISTICA-APLICADA-PARA-LA-ADMINISTRACION}
\providecommand{\itescadepartment}{Estadistica Aplicada para la Administracion}
\providecommand{\itescateachingfigure}{Docente de la materia por definir}
\providecommand{\itescaprogramlevel}{Maestria}
\providecommand{\itescaprogramperiod}{Periodo academico por definir}
\providecommand{\itescamodality}{Modalidad por definir}
\providecommand{\itescaprofessionalprofile}{La Maestria en Gestion Administrativa exige analizar decisiones organizacionales con criterio, evidencia, pertinencia institucional y enfoque de mejora.}
\providecommand{\itescatransfercontext}{Estadistica Aplicada para la Administracion, contexto profesional y aplicacion disciplinar}
\providecommand{\itescalocation}{Cajeme, Sonora}
\providecommand{\itescalocationname}{Cajeme, Sonora}
\providecommand{\itescabibliography}{bibliografia-itesca-mga,ITESCA/maestria-en-gestion-administrativa/estadistica-aplicada-para-la-administracion-mga/estadistica-aplicada-para-la-administracion}
\providecommand{\itescageneralbib}{bibliografia-itesca-mga.bib}
\providecommand{\itescacoursebib}{estadistica-aplicada-para-la-administracion.bib}

\providecommand{\itescamethodologycore}{Problema -> diagnostico -> evidencia -> decision -> consecuencia -> mejora}
\providecommand{\itescawritingmodel}{Redaccion critica, argumentacion ejecutiva y analisis profundo orientado a gestion}
\providecommand{\itescadidactictechniques}{Estudio de caso, matriz de priorizacion, cuadro comparativo, ensayo critico, debate, portafolio y proyecto integrador}

\def\itescadocumenttitle {Estadistica Aplicada para la Administracion}
\def\itescadocumentsubtitle {Contenedor base de materia}
\def\itescadocumentsubject {Estadistica Aplicada para la Administracion}
\def\itescaactivityname {Punto de entrada canonico de materia}
\def\itescasessionname {Estadistica Aplicada para la Administracion}
\def\itescablockname {tronco comun MGA}
\def\itescamainproduct {Informe academico base de la materia}

% ==============================================================================
% PLANTILLA MAESTRA GENERAL - ITESCA
% ============================================================================== 
% Uso:
% 1. Duplica el archivo contenedor de la materia o actividad y actualiza sus
%    metadatos editables.
% 2. Conserva el enfoque de trabajo: no pegar instrucciones completas dentro del
%    documento final; usar la planeacion como lista de cumplimiento.
% 3. Construye productos visuales en LaTeX/TikZ cuando la actividad lo pida.
% 4. Mantiene citas autor-anio con natbib y bibliografia en BibTeX.
% 5. Cierra con criterio propio, no con un resumen plano.
%
% Formato institucional recomendado para ITESCA:
% - Fuente sans serif tipo Helvetica.
% - Interlineado 1.5.
% - Portada institucional, abstract, desarrollo, postura y conclusion.
% - Citas autor-anio con natbib: usar \citep{clave} o \citet{clave}.
% - Bibliografia en BibTeX, consistente con el producto real.
% - Comentarios operativos dentro del archivo para sostener la redaccion.
% - Identidad institucional visible mediante logotipo oficial, marca de agua y
%   vinculacion con la carrera o el contexto profesional.
%
% Criterios generales de cumplimiento:
% - Responder exactamente a la actividad solicitada por la materia o docente.
% - Convertir la consigna en problema, analisis, evidencia y consecuencia.
% - Relacionar el tema con la carrera, la materia y el contexto institucional.
% - Analizar, interpretar y tomar postura; no solo recopilar definiciones.
% - Mostrar producto visible cuando la actividad pida cuadro, esquema, tabla,
%   mapa, diagrama, evidencia o procedimiento.
% - Cerrar con criterio propio, aprendizaje y aplicacion transferible.
%
% Regla para productos visuales:
% - Si la actividad pide cuadro comparativo, mapa, linea de tiempo, matriz,
%   tabla, diagrama o evidencia analitica, construirlo preferentemente en
%   LaTeX/TikZ para mantener nitidez, control y edicion futura.
% - Si la actividad exige herramienta externa, insertar una imagen nitida y
%   conservar la explicacion dentro del documento.
%
% Criterios de redaccion personal:
% - No escribir para llenar espacio; escribir para sostener una idea.
% - Convertir informacion en criterio: que significa, como funciona, que limite
%   presenta, que riesgo aparece y que consecuencia produce.
% - Formula util: concepto -> funcion -> problema -> consecuencia -> postura.
% - La postura personal no debe ser opinion suelta. Debe tener posicion, razon y efecto.
% - La conclusion no debe repetir el desarrollo: debe sintetizar y proyectar.
%
% Control editorial ITESCA:
% - Cada actividad debe tener una ficha de lectura: consigna, producto,
%   evidencia, criterio de evaluacion y transferencia profesional.
% - La identidad institucional no debe limitarse al logotipo. Debe aparecer en
%   la forma de argumentar: precision tecnica, pertinencia academica,
%   vinculacion con carrera y utilidad profesional.
% - Evitar la respuesta generica. El documento debe mostrar que problema atiende
%   la actividad, que decision metodologica se tomo y que evidencia permite
%   sostener el resultado.
% - Cuando la actividad sea de induccion, diagnostico o primer ingreso, traducir
%   el contenido a habitos universitarios: lectura, organizacion, evidencia,
%   integridad academica, autonomia y seguimiento institucional.
% - Cuando la actividad sea disciplinar, traducir el contenido a un proceso,
%   sistema, caso, herramienta, decision administrativa, tecnica o profesional.
%
% Regla de nivel minimo:
% - Si la actividad solo pide identificar o describir, elevar el producto a
%   comprension aplicada: explicar funcion, limite y ejemplo.
% - Si pide cuadro, mapa, matriz, evidencia o conclusion, llegar por lo menos a
%   analisis: comparar criterios, justificar diferencias y cerrar con postura.
% - Si pide propuesta, procedimiento o solucion, llegar a evaluacion/creacion:
%   definir criterios, construir producto propio y justificar su pertinencia.
%
% Voz editorial esperada:
% - Academica, directa, institucional y funcional.
% - No sonar como copia de fuentes ni como resumen escolar.
% - Usar criterio personal con moderacion: posicion, razon y efecto.
% - Antes de redactar, decidir la funcion de cada seccion: abrir problema,
%   delimitar criterio, demostrar con evidencia, valorar o concluir.
%
% Nivel cognitivo esperado:
% - Comprender: explicar el tema con lenguaje propio.
% - Aplicar: llevar el concepto a un caso, proceso, procedimiento o producto.
% - Analizar: distinguir componentes, tensiones, limites y decisiones.
% - Evaluar: sostener un criterio con razones y consecuencias.
% - Crear: construir un producto propio cuando la actividad lo solicite.
%
% Checklist antes de entregar:
% [ ] Titulo, actividad, asignatura, codigo, docente y fecha actualizados.
% [ ] Tecnica didactica identificada y justificada.
% [ ] Existe una tesis, postura o criterio central.
% [ ] La evidencia permite sostener las conclusiones.
% [ ] Se demuestra transferencia academica o profesional.
% [ ] El objetivo coincide con la consigna real.
% [ ] El producto solicitado aparece visible y completo.
% [ ] La introduccion plantea un problema y no solo anuncia el tema.
% [ ] El desarrollo explica, conecta y analiza.
% [ ] El texto aterriza el tema a la carrera o al contexto profesional.
% [ ] Hay citas autor-anio y bibliografia consistente.
% [ ] La conclusion responde que se comprendio y que criterio se sostiene.
% [ ] Quitar o comentar \pendiente cuando el documento final este terminado.
% ============================================================================== 

% ------------------------------------------------------------------------------
% METADATOS EDITABLES
% ------------------------------------------------------------------------------
\providecommand{\itescadocumenttitle}{Plantilla institucional ITESCA}
\providecommand{\itescadocumentsubtitle}{Actividad X}
\providecommand{\itescadocumentsubject}{Informe academico}
\providecommand{\itescastudentname}{Martin Jonathan de la Cruz}
\providecommand{\itescastudentid}{26130503}
\providecommand{\itescacoursename}{Nombre de la asignatura}
\providecommand{\itescacoursecode}{ITESCA}
\providecommand{\itescafaculty}{Programa academico ITESCA}
\providecommand{\itescadepartment}{Departamento academico}
\providecommand{\itescateachingfigure}{Nombre por definir}
\providecommand{\itescastudentemail}{26130503@cajeme.tecnm.mx}
\providecommand{\itescaprogramlevel}{Nivel academico por definir}
\providecommand{\itescaprogramperiod}{Periodo academico por definir}
\providecommand{\itescamodality}{Modalidad por definir}
\providecommand{\itescaactivityname}{Actividad por definir}
\providecommand{\itescasessionname}{Sesion por definir}
\providecommand{\itescablockname}{Bloque por definir}
\providecommand{\itescamainproduct}{Producto academico por definir}
\providecommand{\itescaprofessionalprofile}{Perfil profesional por definir}
\providecommand{\itescainstitutionalaxis}{Formacion academica, identidad institucional y transferencia profesional}
\providecommand{\itescaevidencefocus}{Evidencia academica, producto visible y criterio propio}
\providecommand{\itescaqualitycriterion}{Claridad, pertinencia, evidencia, analisis y cierre transferible}
\providecommand{\itescatransfercontext}{Carrera, asignatura, entorno profesional o trayectoria formativa}
\providecommand{\itescalocation}{Cajeme, Sonora}
\providecommand{\itescabibliography}{bibliografia-itesca-mga,ITESCA/maestria-en-gestion-administrativa/estadistica-aplicada-para-la-administracion-mga/estadistica-aplicada-para-la-administracion}
\providecommand{\itescaasset}{ITESCA/assets-itesca/itesca-monograma.png}
\providecommand{\itescaofficialasset}{ITESCA/assets-itesca/web/logo-itesca-contorno.png}
\providecommand{\itescasecondaryasset}{ITESCA/assets-itesca/itesca-monograma.png}
\providecommand{\itescacampusurl}{https://www.itesca.edu.mx/}
\providecommand{\itescaportalurl}{https://www.itesca.edu.mx/portalacademico/}
\providecommand{\itescavirtualurl}{https://www.itesca.edu.mx/modulos/oferta/ead/index.php}

% CREACION DEL DOCUMENTO
\documentclass[
	spanish,
	letterpaper, oneside
]{article}

% INFORMACION DEL DOCUMENTO
\def\documenttitle {\itescadocumenttitle}
\def\documentsubtitle {\itescadocumentsubtitle}
\def\documentsubject {\itescadocumentsubject}

\def\documentauthor {\itescastudentname}
\def\coursename {\itescacoursename}
\def\coursecode {\itescacoursecode}

\def\universityname {Instituto Tecnologico Superior de Cajeme}
\def\universityfaculty {\itescafaculty}
\def\universitydepartment {\itescadepartment}
\def\universitydepartmentimage {\itescaofficialasset}
\def\universitydepartmentimagecfg {width=5.2cm}
\def\universitylocation {\itescalocation}

% INTEGRANTES, PROFESORES Y FECHAS
\def\authortable {
	\begin{tabular}{ll}
		\textbf{Alumno:}
		& \begin{tabular}[t]{l}
			 \itescastudentname \\
		\end{tabular} \\
		Matricula: & \itescastudentid \\
		Correo institucional: & \itescastudentemail \\
		Asignatura: & \itescacoursename \\
		Codigo: & \itescacoursecode \\
		Nivel: & \itescaprogramlevel \\
		Periodo: & \itescaprogramperiod \\
		Modalidad: & \itescamodality \\

		\textit{Figura docente:}
		& \begin{tabular}[t]{l}
			\itescateachingfigure
		\end{tabular} \\
		Sesion/Bloque: & \itescasessionname\ / \itescablockname \\
		Producto: & \itescamainproduct \\
		& \\
		\multicolumn{2}{l}{Fecha de realizacion: \today} \\
		\multicolumn{2}{l}{\universitylocation}
	\end{tabular}
}

% IMPORTACION DEL TEMPLATE BASE
\input{template}

% FORMATO GENERAL
\usepackage{helvet}
\renewcommand{\familydefault}{\sfdefault}

% IDENTIDAD VISUAL ITESCA
\definecolor{itescaRedDark}{HTML}{8F1117}
\definecolor{itescaRed}{HTML}{D70F18}
\definecolor{itescaCharcoal}{HTML}{141414}
\definecolor{itescaSilver}{HTML}{D9D9D9}
\definecolor{itescaPaper}{HTML}{F7F7F7}
\definecolor{itescaInk}{HTML}{181818}
\definecolor{itescaMuted}{HTML}{6A6A6A}
\colorlet{itescaGreenDark}{itescaRedDark}
\colorlet{itescaGreen}{itescaCharcoal}
\colorlet{itescaGold}{itescaSilver}
\def\portraittitlecolor {itescaRedDark}

% TABLAS Y PRODUCTOS VISUALES
\usepackage{array}
\usepackage{longtable}
\usepackage{pdflscape}
\usepackage{tikz}
\usetikzlibrary{
	arrows.meta,
	positioning,
	calc,
	matrix,
	fit,
	shapes.geometric,
	shadows.blur
}

% CITAS EN FORMATO AUTOR-ANIO, COMPATIBLE CON APA 7 EN EL CUERPO DEL TEXTO
\setcitestyle{authoryear,open={(},close={)}}

% AYUDAS DE REDACCION
% Quitar o comentar \pendiente cuando el documento final este terminado.
\newcommand{\pendiente}[1]{\textcolor{red}{Refuerzo Ciclo A materializado}}

% MARCA DE AGUA EN PORTADA
\def\coverwatermarkenabled {true}
\def\coverwatermarkimage {\itescaasset}
\def\coverwatermarkopacity {0.13}
\def\coverwatermarkwidth {0.58\paperwidth}
\def\coverwatermarkxshift {0cm}
\def\coverwatermarkyshift {-0.10cm}

\newcommand{\insertcoverwatermark}{%
	\ifthenelse{\equal{\coverwatermarkenabled}{true}}{%
		\AddToShipoutPictureBG*{%
			\begin{tikzpicture}[remember picture,overlay]
				\node[
					opacity=\coverwatermarkopacity,
					inner sep=0pt,
					xshift=\coverwatermarkxshift,
					yshift=\coverwatermarkyshift
				] at (current page.center) {%
					\includegraphics[width=\coverwatermarkwidth]{\coverwatermarkimage}%
				};
			\end{tikzpicture}%
		}%
	}{}%
}

\newcommand{\insertcoverinstitutionalfooter}{%
	\AddToShipoutPictureFG*{%
		\begin{tikzpicture}[remember picture,overlay]
			\fill[itescaCharcoal] (current page.south west) rectangle ([yshift=1.35cm]current page.south east);
			\fill[itescaRed] ([yshift=1.35cm]current page.south west) rectangle ([yshift=1.47cm]current page.south east);
			\node[anchor=south west] at ([xshift=0.95cm,yshift=0.16cm]current page.south west) {%
				\includegraphics[height=0.88cm]{\itescasecondaryasset}%
			};
			\node[
				anchor=south east,
				text=white,
				align=right,
				font=\sffamily\bfseries\small
			] at ([xshift=-0.95cm,yshift=0.28cm]current page.south east) {%
				Instituto Tecnologico Superior de Cajeme\\
				\textcolor{itescaSilver}{\itescalocation \textperiodcentered\ \href{\itescacampusurl}{itesca.edu.mx}}%
			};
		\end{tikzpicture}%
	}%
}

\usepackage{pdfpages}
\usepackage{ragged2e}
% Preliminares locales de Estadística. Se cargan desde el contenedor de materia.
% Resumen e índice son aparato preliminar: no añaden actos al cuerpo académico.
\providecommand{\estadisticadeclaracionia}{%
Se utilizó inteligencia artificial como apoyo para la corrección y organización del
contenido. El autor declara haber realizado y validado la investigación documental,
el análisis y los resultados presentados, y asume su responsabilidad académica.%
}
\providecommand{\estadisticapreliminares}{%
  \ifdefined\estadisticaresumen\else
    \errmessage{Definir estadisticaresumen con la sintesis especifica del informe}%
  \fi
  \section*{Resumen}
  \estadisticaresumen\footnote{\textbf{Declaración de uso de IA.} \estadisticadeclaracionia}
  \par\medskip
  \begingroup
    \setcounter{tocdepth}{2}%
    \renewcommand{\contentsname}{Tabla de contenido}%
    \tableofcontents
  \endgroup
  \clearpage
}

\def\estadisticaresumen{La evidencia administrativa necesita una cadena coherente entre problema, población, datos y método. El informe articula selección de muestras, estimación, pruebas de hipótesis y relaciones estadísticas, con ejemplos de satisfacción, facturación y planificación. Explica por qué una cifra precisa puede seguir sesgada, por qué la significación no equivale a importancia práctica y por qué una asociación no prueba causalidad. Se sostiene que los resultados deben comunicar incertidumbre y límites, junto con criterios para revisar las decisiones.}
\begin{document}
% Portada histórica preservada literalmente, incluida su fecha y metadatos originales.
\includepdf[pages=1,pagecommand={}]{ITESCA/maestria-en-gestion-administrativa/estadistica-aplicada-para-la-administracion-mga/assets-estadistica-aplicada-para-la-administracion/portada-historica-mga.pdf}
\templatePagecfg
\onehalfspacing
\templateFinalcfg
% Etiquetas breves en el pie para evitar solapamiento del código histórico extenso.
\fancyfoot[L]{\small\textit{Datos e incertidumbre}}
\fancyfoot[R]{\small\textit{GGEA01 / Estadística aplicada}}
\hypersetup{hidelinks}
\justifying
\estadisticapreliminares
% Contrato: tres actos. Método didáctico: recorrido conceptual con contraste aplicado.
\section{Introducción}
Las decisiones de compra, contratación y asignación de recursos se toman con información incompleta. La estadística aplicada permite ordenar registros, inferir características de una población y evaluar la incertidumbre de esas inferencias. La dificultad no consiste únicamente en calcular una media o un porcentaje: exige saber a quién representan los datos, qué pregunta responden y qué condiciones permiten trasladar el resultado a una decisión \citep{holmesBusiness2023}.

Un análisis administrativo defendible vincula problema, diseño, evidencia y consecuencia. La precisión de una cifra no compensa la exclusión sistemática de clientes ni convierte una asociación en una causa. El criterio que articula los métodos es elegirlos por la naturaleza de los datos y por el riesgo de equivocarse, manteniendo visible el alcance de cada conclusión.

\section{Datos, inferencia y decisiones de gestión}
\subsection{Del problema administrativo a la población}
Una pregunta sobre satisfacción debe delimitar clientes, periodo y canal; una pregunta sobre errores contables debe definir operaciones y criterio de error. La población objetivo expresa ese alcance y el marco muestral permite acceder a sus unidades. Si el registro solo contiene ventas digitales, una selección aleatoria de ese registro no representa automáticamente a los compradores de tienda. El mecanismo de selección y las respuestas efectivas condicionan la interpretación del estadístico \citep{holmesBusiness2023,bradley2021}.

El muestreo probabilístico asigna probabilidades de inclusión conocidas y positivas dentro de la población cubierta. El estratificado permite incluir grupos pertinentes; el muestreo por conglomerados puede reducir costos de acceso. Una encuesta por conveniencia puede ser útil para explorar necesidades, pero no hereda sin justificación el margen de error de una muestra aleatoria. La cobertura y la no respuesta se examinan junto al tamaño: más observaciones no eliminan por sí solas un sesgo de participación \citep{bradley2021}.

\subsection{Estimar y comunicar la incertidumbre}
Un parámetro describe una población y un estimador es una regla calculada a partir de la muestra. La media muestral cambia entre selecciones; su distribución muestral describe esa variación. Para observaciones independientes e idénticamente distribuidas con varianza finita, su varianza es la poblacional dividida entre el tamaño de muestra. El error estándar expresa la dispersión del promedio y no debe confundirse con la desviación de observaciones individuales \citep{wackerly2010}.

La estimación puntual facilita presupuestar con un escenario central. Un intervalo de confianza añade límites bajo un procedimiento y supuestos definidos. La interpretación frecuentista del 95\% se refiere a la cobertura del procedimiento en muestreos repetidos; no asegura que el 95\% de las ventas futuras quede entre los extremos ni asigna esa probabilidad a un parámetro fijo. La amplitud depende de confianza, variabilidad, tamaño y diseño \citep{nistHandbook}.

Para una encuesta hipotética de satisfacción, 95\% de confianza, margen de cinco puntos porcentuales y proporción de planificación 0.5 producen 384.16 observaciones mediante la aproximación normal para una población grande; se redondea a 385. Con respuesta prevista de 80\%, se planifican 482 invitaciones. El cálculo supone selección aleatoria simple y no garantiza la respuesta ni corrige exclusiones del marco. La cifra adquiere significado cuando se documentan esas condiciones \citep{holmesBusiness2023}.

\subsection{Contrastar hipótesis con relevancia práctica}
Un contraste exige definir parámetro, hipótesis, nivel de significación, diseño y estadístico. Un resultado significativo puede acompañar una diferencia pequeña para la operación; un resultado no significativo puede provenir de información insuficiente. No rechazar la hipótesis nula no demuestra su verdad. Una decisión de cambiar un proceso debe considerar magnitud del efecto, incertidumbre, costos y supuestos, además del resultado del contraste \citep{nistHandbook}.

La comparación de una o dos muestras depende de si las observaciones son independientes o pareadas y de la escala de medición. Los procedimientos no paramétricos, incluidos los basados en rangos, ofrecen alternativas bajo condiciones concretas. No eliminan la necesidad de independencia, una selección adecuada ni una interpretación cuidadosa de las distribuciones. Escoger la prueba después de buscar repetidamente un resultado favorable debilita el fundamento de la decisión.

\subsection{Relaciones, predicción y límites del modelo}
Una relación estadística entre inversión comercial y ventas no prueba por sí sola que el gasto adicional cause el aumento. Pueden intervenir temporada, precios, demanda y selección de mercados. La regresión describe una relación condicionada al modelo y a los datos disponibles; sus supuestos y residuos deben revisarse antes de extrapolar. Del mismo modo, un intervalo para una media y un intervalo para una observación futura responden a preguntas diferentes \citep{nistHandbook}.

La planificación debe distinguir entre describir lo observado y predecir lo que ocurrirá. Un modelo útil para explicar variabilidad pasada necesita validación con evidencia pertinente para respaldar compromisos de inventario o contratación. Los cambios del entorno pueden invalidar relaciones históricas. Mantener seguimiento y comparar resultados con previsiones permite ajustar una decisión sin tratar el ajuste original como garantía permanente.

La integración de estos métodos conduce a una conclusión de gestión: la calidad del análisis reside en la cadena completa de inferencia. Un promedio bien calculado no corrige un marco incompleto; una prueba significativa no identifica automáticamente causalidad; y una predicción precisa en registros pasados puede fallar si cambian las condiciones. Por ello, solicitar población, diseño, supuestos, magnitud e incertidumbre resulta más útil que pedir únicamente un indicador final \citep{holmesBusiness2023,nistHandbook,bradley2021}.

\section{Conclusiones}
La estadística aplicada transforma información parcial en decisiones revisables cuando mantiene correspondencia entre población, medición y método. Muestreo, estimación, contraste y modelación cumplen funciones distintas y complementarias. Su aplicación requiere comunicar qué se conoce, bajo qué condiciones y qué incertidumbre permanece.

Considero que una gerencia debe exigir resultados acompañados de su alcance y de las razones que justifican el procedimiento. La razón es que las consecuencias económicas de decidir con una cifra mal interpretada pueden superar el costo de recabar información adicional. La consecuencia profesional es formular umbrales de decisión, reconocer límites de inferencia y establecer seguimiento para revisar acciones ante evidencia nueva.

El aprendizaje central consiste en usar el cálculo como parte de un argumento y no como sustituto del juicio. Una decisión responsable no promete certeza: demuestra por qué la evidencia disponible justifica una acción proporcional al riesgo y deja explícito qué podría hacerla cambiar.
\clearpage
\renewcommand{\bibsection}{\section*{Referencias}\addcontentsline{toc}{section}{Referencias}}
\setlength{\bibhang}{1.27cm}\setlength{\bibsep}{10pt}
\begin{thebibliography}{4}
\bibitem[Bradley et al.(2021)]{bradley2021}
Bradley, V. C., Kuriwaki, S., Isakov, M., Sejdinovic, D., Meng, X.-L., \& Flaxman, S. (2021). Unrepresentative big surveys significantly overestimated US vaccine uptake. \textit{Nature, 600}(7890), 695--700. \url{https://doi.org/10.1038/s41586-021-04198-4}
\bibitem[Holmes et al.(2023)]{holmesBusiness2023}
Holmes, A., Illowsky, B., \& Dean, S. (2023). \textit{Introductory business statistics} (2.\textsuperscript{a} ed.). OpenStax. \url{https://openstax.org/details/books/introductory-business-statistics-2e}
\bibitem[NIST y SEMATECH(2003)]{nistHandbook}
National Institute of Standards and Technology, \& SEMATECH. (2003). \textit{NIST/SEMATECH e-Handbook of statistical methods}. \url{https://doi.org/10.18434/M32189}
\bibitem[Wackerly et al.(2010)]{wackerly2010}
Wackerly, D. D., Mendenhall, W., \& Scheaffer, R. L. (2010). \textit{Estadística matemática con aplicaciones} (7.\textsuperscript{a} ed.). Cengage Learning Editores.
\end{thebibliography}

\end{document}
